\section{Independent verification of
WHR3--WHR7}\label{independent-verification-of-whr3whr7}

25 September 2026. This private verification concerns only the named
receiving-space argument. It assigns no parity, addition, coefficient,
or Hilbert structure to primitive source \texttt{Z1/tau}. The reviewed
manuscript was not edited.

\subsection{Task record and reading
coverage}\label{task-record-and-reading-coverage}

This independent mathematical verification checks WHR3--WHR7: the
explicit even Schwartz generator, Fourier multiplier annihilator in the
weighted dual space, derivative integrability cutoff, convergence of
frequency cutoffs, and exact closure and Gram model. All Fourier factors
and signs are retained. Primitive Z\_1/tau has no parity or addition;
these operations occur in the specified receivers.

Read the complete WHR manuscript. Also read the actual RZ5 formulas in
\href{https://github.com/KokunoYumeto/zeta-function-research-reader/blob/36a82ecca98addb8a98b48204e349737856c7223/workbenches/splitzero-tandem/continuations/20260924-cohomology-weight-and-character-lifts/gct-weight-control/proofs/ORIGINAL_ZETA_RESOLVENT_AND_PROJECTORS.md}{sources/programme\_private/source\_comparison/ORIGINAL\_ZETA\_RESOLVENT\_AND\_PROJECTORS.md},
together with its definitions of the shifted Mellin transform and the
function space. Read original author TeX
\href{https://arxiv.org/abs/math/9811068v1}{../connes98\_intake/math\_9811068v1\_author.tex},
source lines 791--931 and 3795--3946, to verify the stated exponent
comparison and multiplicity cutoff. This check treats the SSI equality
\texttt{M\_0\ J\ =\ I} as the named input theorem, rather than claiming
to reprove SSI. No numerical zero sample or zero-spacing assumption is
used.

\subsection{V1. Generator and original Mellin
factors}\label{v1.-generator-and-original-mellin-factors}

Put \[
p(z)=\frac1{2\sqrt\pi}e^{-z^2/4}.
\] On the positive half-line the proposed generator is exactly \[
h_0(v)=\frac14(\log v-2)p(\log v).
\] Every order-\(k\) derivative is a finite sum of
\(v^{-k}P(\log v)e^{-(\log v)^2/4}\). After multiplication by \(v^N\),
the substitution \(z=\log v\) gives a polynomial times
\(\exp(-z^2/4+(N-k)z)\), which tends to zero as \(z\to\pm\infty\). In
particular every derivative tends to zero at the origin, the extension
through the origin is smooth and flat, and the reflected definition is
even and Schwartz.

Gaussian integration and differentiation give \[
\int_{\mathbb R}p(z)e^{sz}\,dz=e^{s^2},\qquad
\int_{\mathbb R}zp(z)e^{sz}\,dz=2s e^{s^2}.
\] Thus the exact Mellin formula is \[
\mathcal M_S h_0(s)=\frac14(2s-2)e^{s^2}
=\frac{s-1}{2}e^{s^2}.
\] Evenness gives \(\int_{\mathbb R}h_0=2\mathcal M_Sh_0(1)=0\), so
\(h_0\in S\). For \(\Re s>1\), absolute convergence permits the
substitution \(v=nu\) in every summand: \[
\mathcal M_0(\Sigma h_0)(s)
=2\sum_{n\ge1}n^{-s}\mathcal M_Sh_0(s)
=(s-1)\zeta(s)e^{s^2}.
\] The factor two in \(\Sigma\) cancels precisely the displayed
denominator two in \(\mathcal M_Sh_0\). The equality extends to all
\(s\), retaining the removable value \(\exp(1)\), the real exponential
value, at \(s=1\). The factors \(s-1\) and \(e^{s^2}\) have no zeros on
\(\Re s=1/2\), so each zero on that line has exactly its original zeta
multiplicity. The trivial zeros still occur in the entire Mellin
formula.

The standard source summation result already used in WHR1 puts
\(b_0=\Sigma h_0\) in \(A\). This membership can also be checked
directly: for large \(u\), differentiate the original sum and use
Schwartz estimates; for small \(u\), Poisson summation, \(h_0(0)=0\),
and \(\widehat h_0(0)=0\) give \[
b_0(u)=u^{-1}\sum_{n\ne0}\widehat h_0(n/u),
\] whose differentiated series has arbitrarily rapid decay at zero.
Consequently \(g_0(x)=e^{x/2}b_0(e^x)\) is Schwartz and \[
\widehat g_0(t)=(-\tfrac12+it)\zeta(\tfrac12+it)e^{(1/2+it)^2}.
\]

\subsection{V2. Weighted dual and cutoff
topology}\label{v2.-weighted-dual-and-cutoff-topology}

Write \(K_{-\delta}=L^2((1+x^2)^{-\delta}dx)\), and use the bilinear
dual pairing \(\ell_y(g)=\int gy\). Weighted Cauchy--Schwarz identifies
its exact norm with \(\|y\|_{K_{-\delta}}\); surjectivity follows from
the Hilbert representation theorem after conjugating the representative
and multiplying by the weight. Every such \(y\) is tempered because
Schwartz functions belong continuously to \(K_\delta\).

For translation \(\tau_z y(x)=y(x-z)\), changing variables and using \[
1+t^2\leq2(1+(t+z)^2)(1+z^2)
\] give \[
\|\tau_z y\|_{K_{-\delta}}
\leq2^{\delta/2}(1+z^2)^{\delta/2}\|y\|_{K_{-\delta}}.
\] This is the stated bound, with the correct direction of the negative
weight. Smooth compactly supported functions are dense by truncation and
smoothing, because the weight and its reciprocal are bounded on compact
intervals. The displayed bound and dominated convergence on compactly
supported functions therefore imply strong continuity of all
translations.

The frequency distribution in WHR4 is \[
T_y=\frac1{2\pi}\mathcal F_-y,
\qquad \mathcal F_-f(t)=\int f(x)e^{-itx}\,dx.
\] Its inverse is \(y(x)=\langle T_y(t),e^{itx}\rangle\), interpreted in
distributions. Accordingly multiplication by \(\chi_R(t)=\chi(t/R)\)
gives \[
y_R=k_R*y,\qquad
k_R(x)=Rk(Rx),\qquad
k(x)=\frac1{2\pi}\int\chi(t)e^{itx}\,dt.
\] There is no sign reversal in this kernel. Fourier inversion gives
\(\int k=\chi(0)=1\). The exact uniform estimate needed for cutoff
convergence is \[
\|k_R*y\|_{K_{-\delta}}
\leq 2^{\delta/2}\|y\|_{K_{-\delta}}
\int_{\mathbb R}|k(z)|(1+(z/R)^2)^{\delta/2}\,dz
\leq 2^{\delta/2}\|y\|_{K_{-\delta}}
\int_{\mathbb R}|k(z)|(1+z^2)^{\delta/2}\,dz.
\] The last integral is finite because \(k\) is Schwartz. Moreover \[
y_R-y=\int k(z)(\tau_{z/R}y-y)\,dz.
\] The norm of the integrand is bounded by \[
|k(z)|\bigl(2^{\delta/2}(1+z^2)^{\delta/2}+1\bigr)
\|y\|_{K_{-\delta}},
\] an integrable function independent of \(R\ge1\), and tends to zero
for every fixed \(z\). This proves convergence in the actual weighted
dual norm. The same convolution estimate proves boundedness of every
fixed smooth compactly supported frequency multiplier.

\subsection{V3. Multiplier equation and every point
derivative}\label{v3.-multiplier-equation-and-every-point-derivative}

The dilation formula is \[
UT_{e^v}b_0(x)=e^{v/2}g_0(x-v).
\] If \(\ell_y\) annihilates \(J\), then \(C(v)=\int y(x)g_0(x-v)dx=0\)
for all real \(v\). Here \(C=y*\check g_0\), with
\(\check g_0(x)=g_0(-x)\). Consequently \[
\mathcal F_-C
=(\mathcal F_-y)(\mathcal F_-\check g_0)
=(2\pi T_y)m,
\qquad
\frac1{2\pi}\mathcal F_-C=mT_y=0.
\] This calculation explicitly verifies both Fourier factors. The
convolution identity is valid in tempered distributions: weighted
translation bounds make the displayed convolution a smooth function of
polynomial growth, and testing it against Schwartz functions justifies
the convolution theorem. Equivalently, for every Schwartz \(\psi\), \[
\int C(v)\psi(v)\,dv
=\langle T_y,m\widehat\psi\rangle.
\] Since the positive-phase Fourier transform maps Schwartz space onto
itself, this also proves \(mT_y=0\) directly.

Outside the zeros of \(m\), division of compactly supported test
functions by \(m\) proves \(T_y=0\). In a neighborhood of a zero
\(\gamma\) of order \(q=m_\rho\), write \(m(t)=(t-\gamma)^qv(t)\) with
\(v\) nonzero. Multiplication by the smooth local inverse of \(v\) gives
\((t-\gamma)^qT_y=0\). Choose a compactly supported cutoff equal to one
on a smaller neighborhood of \(\gamma\). Subtracting this cutoff times
the degree-\((q-1)\) Taylor polynomial of an arbitrary test function
leaves a multiple of \((t-\gamma)^q\) near the support of the localized
distribution. The localized distribution therefore depends only on the
derivatives of orders \(0,\ldots,q-1\), and is a linear combination of
precisely the point derivatives of those orders.

Choose the local cutoff equal to one near \(\gamma\), with no other zero
in its support. V2 ensures that the corresponding function still belongs
to \(K_{-\delta}\). Explicit inversion of the point derivatives gives \[
T=\sum_{j=0}^{q-1}a_j\delta_\gamma^{(j)}
\quad\Longleftrightarrow\quad
y(x)=e^{i\gamma x}\sum_{j=0}^{q-1}a_j(-ix)^j.
\] If the nonzero polynomial on the right has degree \(d\), then its
absolute square is bounded above and below by positive multiples of
\(|x|^{2d}\) for all sufficiently large \(|x|\). Thus its weighted
square integral converges exactly when \(2d-2\delta<-1\), equivalently
\(d<\delta-1/2\). At equality it diverges logarithmically. This proves
the strict point-derivative cutoff without permitting a cancellation
between different zeros.

For the original derivative functional, the point-distribution factor is
\[
T_{\rho,j}=i^j\delta_{\gamma_\rho}^{(j)}.
\] Indeed \[
\langle i^j\delta_{\gamma}^{(j)},\widehat g\rangle
=i^j(-1)^j\int (ix)^jg(x)e^{i\gamma x}\,dx
=\int x^jg(x)e^{i\gamma x}\,dx.
\] Both factors \(i^j\) are needed, and their product cancels \((-1)^j\)
exactly.

Finally \(\chi_RT_y\) has finite support, because the zeros of the
nonzero analytic \(m\) are locally finite. It is a finite linear
combination of the allowed point derivatives. V2 proves that these
finite combinations converge in the weighted dual norm. This is the
global converse in WHR5.2; no separation estimate on distinct zero
ordinates is needed. The forward inclusion follows from continuity of
the original derivatives and their vanishing on every actual
\(\Sigma h\). Therefore WHR5.2 is correct globally.

\subsection{V4. Closure, derivative count, and the complete sequence
space}\label{v4.-closure-derivative-count-and-the-complete-sequence-space}

For a closed subspace of a Hilbert space, separation by the nonzero
orthogonal component proves that the common kernel of its continuous
annihilator is exactly that subspace. Applying this fact to the
norm-closed span just proved gives \[
\overline J^{H_\delta}
=\bigcap_{(\rho,j)\in\mathscr J_\delta}\ker L_{\rho,j}.
\] The number of nonnegative integers satisfying \(j<\delta-1/2\),
capped by \(m_\rho\), is exactly \[
\min\{m_\rho,\max(0,\lceil\delta-1/2\rceil)\}.
\] At a half-integer \(\delta=k+1/2\), only \(j=0,\ldots,k-1\) are
included. The count agrees with Connes's original largest integer
\(n<(1+\delta_{\rm CC})/2\) after \(\delta_{\rm CC}=2\delta\): its upper
threshold is \(\delta+1/2\), and the largest integer strictly below that
threshold is \(\lceil\delta-1/2\rceil\). The Connes theorem is stated
for \(\delta_{\rm CC}>1\), corresponding to the nonzero-jet range here.

For finitely supported \(c\), the dual element
\(\sum c_{\rho,j}L_{\rho,j}\) has exactly the norm \(D_\delta(c)\)
displayed in WHR6.4. This norm is nonzero unless \(c=0\). To verify the
needed independence, write a finite sum as
\(\sum_\gamma e^{i\gamma x}P_\gamma(x)\). For a chosen \(\gamma_0\),
apply \[
\prod_{\gamma\ne\gamma_0}(\partial_x-i\gamma)^{d_\gamma+1},
\qquad d_\gamma=\deg P_\gamma.
\] It kills every summand with \(\gamma\ne\gamma_0\). On the remaining
summand it acts as \(e^{i\gamma_0x}\) times \[
\prod_{\gamma\ne\gamma_0}
(\partial_x+i(\gamma_0-\gamma))^{d_\gamma+1}P_{\gamma_0}.
\] Each factor is invertible on the finite-dimensional polynomial space
of degree at most \(\deg P_{\gamma_0}\), because it is a nonzero scalar
times the identity plus a nilpotent operator. Hence a vanishing initial
sum forces \(P_{\gamma_0}=0\). Applying this to every frequency proves
independence.

Let \(E\) denote the annihilator, with the dual norm. Every sequence
\(v\) with finite \(N_\delta(v)\) defines a bounded complex-linear
functional \[
\sum c_{\rho,j}L_{\rho,j}\longmapsto
\sum c_{\rho,j}v_{\rho,j}
\] on the dense finite span in \(E\). It extends uniquely to \(E\).
Since \(E=(H_\delta/J_\delta)'\) is Hilbert, reflexivity identifies
\(E'\) isometrically with \(H_\delta/J_\delta\), through the canonical
evaluation map. This proves the claimed surjectivity of the sequence
description and its exact norm. No complex conjugate is missing from the
numerator of WHR6.5: the pairing there is between a complex-linear
functional and its argument, rather than a Hilbert inner product. If the
index set is empty, the sole empty sequence has norm zero, agreeing with
the zero quotient.

\subsection{V5. Gram entries, original quotient map, and exact
kernel}\label{v5.-gram-entries-original-quotient-map-and-exact-kernel}

Expanding the absolute square in \(D_\delta(c)^2\) gives \[
D_\delta(c)^2
=\sum_{(\rho,j),(\sigma,k)}
c_{\rho,j}\overline{c_{\sigma,k}}
\int x^{j+k}e^{i(\gamma_\rho-\gamma_\sigma)x}
(1+x^2)^{-\delta}\,dx.
\] Thus the phase difference and complex conjugations in WHR6.7 are
correct. Allowed indices obey \(j+k<2\delta-1\), so every integral
converges absolutely. For the zeroth derivative, the Gamma integral and
Gaussian Fourier integral give exactly \[
\frac{\sqrt\pi}{\Gamma(\delta)}
\int_0^\infty t^{\delta-3/2}
\exp\left(-t-\frac{(\gamma_\rho-\gamma_\sigma)^2}{4t}\right)dt.
\] The relevant index set is nonempty only if \(\delta>1/2\), which also
justifies absolute Fubini in this calculation. Differentiating with
respect to the frequency difference \(j+k\) times produces
\((ix)^{j+k}\). The factor needed to recover the stated Gram entry is
therefore \(i^{-(j+k)}\), as written. The same strict inequality
\(j+k<2\delta-1\) justifies all these derivatives, including when the
frequency difference is zero.

Using the specified SSI input \(\mathcal M_0J=I\), the common-kernel
formula gives \[
A\cap J_\delta=\mathcal M_0^{-1}I_\delta,
\qquad
\ker(A/J\longrightarrow H_\delta/J_\delta)=I_\delta/I.
\] The image is exactly the allowed derivative sequence of a single
entire \(F\in\mathcal B\), and is dense because \(A\) is dense in
\(H_\delta\). It need not be the full Hilbert sequence space. For fixed
derivative order, a Cauchy circle of fixed radius around
\(1/2+i\gamma\), together with the vertical-strip seminorms of \(F\),
proves its stated rapid decay as \(|\gamma|\to\infty\).

The quotient norm is the infimum over the closure \(J_\delta\). It
equals the infimum over \(J\), because the latter is dense in the former
and the norm is continuous. This proves both equalities in WHR7.5.

RZ5 defines \(E_{\rho,j}=e_\rho(s)(s-\rho)^j\) with
\(e_\rho=1+O((s-\rho)^{m_\rho})\) at \(\rho\), and full required
vanishing at every other zero. Its actual derivative is \(j!\) at the
selected \(j\)-th coordinate, with all other relevant derivatives zero.
Taking the raw inverse Mellin transform exactly as defined in WHR7
therefore preserves this factorial and introduces no factor two. These
jets prove that every off-critical primary block lies in the kernel; on
a critical block the kernel is precisely the span of the basis elements
with \(j\ge r_\delta(\rho)\). This is a statement about each genuine
finite-dimensional primary block, with no assertion that the full
quotient is an unrestricted product of them.

\subsection{Verification result and optional precision
edits}\label{verification-result-and-optional-precision-edits}

WHR3--WHR7 are mathematically valid with their stated SSI input. No
substantive correction is required. All displayed factors and signs
checked above are correct. Two optional clarifications would make the
prose more explicit without changing any formula:

\begin{enumerate}
\def\labelenumi{\arabic{enumi}.}
\tightlist
\item
  In the Fourier-convolution explanation after WHR5.4, write both
  equalities \texttt{F\_-\ C\ =\ (2\ pi\ T\_y)m} and
  \texttt{(2\ pi)\^{}(-1)\ F\_-\ C\ =\ m\ T\_y}. This distinguishes the
  divided transform used for \texttt{T\_y} from the undivided transform
  of \texttt{check\ g\_0}.
\item
  In the local cutoff argument, specify that the isolating cutoff equals
  one in a neighborhood of the selected zero. This ensures that its
  derivative coefficients are the actual local coefficients before
  applying the weighted integrability cutoff.
\end{enumerate}

All verification work is complete for the delegated range. The root
manuscript remains unchanged.

\subsection{Publication source and dependency
links}\label{publication-source-and-dependency-links}

The source geometry is Alain Connes and Caterina Consani, \emph{Schemes
over F1 and zeta functions},
\href{https://arxiv.org/abs/0903.2024v3}{arXiv:0903.2024v3, §5}. The
weight-control comparison is Pierre Deligne, \emph{La conjecture de
Weil. II}, Publications mathematiques de l'IHES52 (1980),137-252,
\href{https://www.numdam.org/item/PMIHES_1980__52__137_0/}{§§3.3.11
and3.6}. Deligne material was read through the identified French
transcription and separately recorded peer source-page checks, not
Deligne-authored TeX. These sources are not asserted to contain the new
programme derivations. The
\href{https://github.com/KokunoYumeto/zeta-function-research-reader/blob/36a82ecca98addb8a98b48204e349737856c7223/workbenches/splitzero-tandem/continuations/20260924-cohomology-weight-and-character-lifts/BUILD_AND_REVIEW.md}{preceding
complete proof and reading record} records inherited source versions and
actual inspection limits. The investigator's corrected construction
remains attributed in the original text above.
